\chapter{Primitive Heisenberg coefficients}
\label{cbc:chap:heisenberg}

Modulo derivatives, a rational differential with poles at one fixed
point is determined by its residue. Moving the points changes this calculation: the
coefficients of partial fractions can acquire poles on the collision
diagonals. The Heisenberg coefficient must retain these terms before
the points are identified. Its primitive cohomology is flat over the
position ring, but its diagonal constituents form nonsplit extensions.
Those extensions determine both the collision maps and their duals.

Throughout this chapter, $k$ is a field of characteristic zero.
Complexes have cohomological grading, and differential modules are
left modules. All rational functions have finite pole order. No
completion in pole order is taken.

\section{Principal parts and differentiation}
\label{cbc:hc:principal-parts}

Let $n\geq1$, let $z_1,\ldots,z_n$ be the position coordinates, and
let $w$ be an auxiliary coordinate. Set
\begin{equation}\label{cbc:hc:poles}
 R_n=k[z_1,\ldots,z_n],\qquad
 P(w)=\prod_{i=1}^n(w-z_i),\qquad
 M_n=R_n[w,P^{-1}]/R_n[w].
\end{equation}

Division by the monic denominator gives each element of $M_n$ a
unique proper rational representative. Differentiation in $w$ and
in each position preserves the polynomial submodule and therefore
acts on $M_n$. These operators commute. The primitive coefficient is
\begin{equation}\label{cbc:hc:primitive}
 H_n=\operatorname{coker}(\partial_w:M_n\longrightarrow M_n).
\end{equation}
The differential $f(w)\,dw$ records the residue convention for a
representative $f(w)$ of this quotient.

For a nonzero level $\kappa\in k$, let $xM_n$ and $\gamma M_n$
be copies of $M_n$ in degrees zero and one. The linear
Heisenberg coefficient complex is
\begin{equation}\label{cbc:hc:linear-complex}
 K_{\kappa,n}=[xM_n\xrightarrow{-\kappa\partial_w}\gamma M_n].
\end{equation}
Its degree-one cohomology is $H_n$. The polynomial coefficient
is its graded symmetric algebra, treated in
Section~\ref{cbc:hc:polynomial}. The primitive calculation
determines its cohomology.

At one point every class of pole order greater than one is a
$w$-derivative, and $H_1$ is free on the class of
$(w-z_1)^{-1}$.
At two points, writing
\[
 g_i=[(w-z_i)^{-1}],\qquad
 g_{12}=[((w-z_1)(w-z_2))^{-1}],
\]
gives
\begin{equation}\label{cbc:hc:binary-attachment}
 (z_1-z_2)g_{12}=g_1-g_2.
\end{equation}
This identity holds over the polynomial position ring. Dividing by
$z_1-z_2$ describes the separated residues, but loses the question
of how they extend over the diagonal. Finite denominator stages
compute the quotient over that ring.

Put $M^{(0)}=0$ and
$M^{(m)}=P^{-m}R_n[w]/R_n[w]$ for $m\geq1$, with arity $n$ fixed. The elements $w^e/P^m$, $0\leq e<mn$, form an
$R_n$-basis of $M^{(m)}$. Repeated monic division also gives a basis
of the whole module:
\begin{equation}\label{cbc:hc:free-basis}
 M_n=\bigoplus_{m\geq1}\bigoplus_{e=0}^{n-1}
           R_n\frac{w^e}{P^m}.
\end{equation}
Indeed, multiplying a finite relation by its largest denominator
and reducing modulo $P$ first kills the numerator at that
denominator. Induction kills every numerator.

\begin{theorem}[Finite pole stages]\label{cbc:hc:stages}
For $m\geq0$, the module
$F_m=M^{(m+1)}/\partial_wM^{(m)}$ is free of rank $n$, with ordered
representatives
\begin{equation}\label{cbc:hc:stage-basis}
 c_{m,0}=\frac1{P^{m+1}},\quad\ldots,\quad
 c_{m,n-2}=\frac{w^{n-2}}{P^{m+1}},\qquad
 c_{m,n-1}=\frac{w^{n(m+1)-1}}{P^{m+1}}.
\end{equation}
For $n=1$ only the last representative occurs. The inclusions of
principal parts induce transitions $T_m:F_m\to F_{m+1}$, and
$H_n=\operatorname*{colim}_mF_m$. In these bases,
\begin{equation}\label{cbc:hc:transition-determinant}
 \det T_m=(-1)^{n(n-1)/2}
 \frac{(m+1)^{n-1}}
 {n\prod_{j=1}^{n-1}(n(m+1)-j)}
 \prod_{i<j}(z_i-z_j)^2.
\end{equation}
At a geometric position with exactly $r$ distinct points,
the specialized transition has rank $r$.
\end{theorem}

\begin{proof}
For $0\leq e<mn$, differentiation has numerator
\[
 \partial_w\frac{w^e}{P^m}
 =\frac{d_{m,e}}{P^{m+1}},\qquad
 d_{m,e}=e w^{e-1}P-mw^eP'.
\]
Its leading term is $(e-mn)w^{e+n-1}$. These nonzero integer
pivots occupy the distinct degrees $n-1,\ldots,n(m+1)-2$.
Descending elimination leaves precisely
\eqref{cbc:hc:stage-basis}. It also proves that
$\partial_w:M_n\to M_n$ is injective. Every fraction and every
primitive belong to finite stages, so the colimit of these
cokernels is $H_n$.

The transition is computed by multiplying a basis numerator by
$P$ and eliminating the pivots for $\partial_wM^{(m+1)}$.
Only nonzero integers are divided out. Its entries are therefore
polynomials in the coefficients of $P$.

After specialization to an algebraically closed field, a proper
rational differential has a proper rational primitive if and only
if its residues vanish. At a root of multiplicity $d$ in $P$, a
zero-residue differential with denominator $P^{m+1}$ has a primitive
of pole order at most $(m+1)d-1$. That primitive lies in $M^{(m+1)}$.
Consequently the kernel of the specialized $T_m$ is exactly the
zero-residue subspace of $F_m$. The simple poles at the $r$
remaining points give independent residue vectors and belong to
$M^{(m+1)}$. This proves the rank assertion.

The determinant vanishes on every pair diagonal. It is symmetric
in the roots, so each difference $z_i-z_j$ divides it to even
order. The last representative in \eqref{cbc:hc:stage-basis} has
total finite residue one, and all the others have total residue
zero. The last row of $T_m$ is thus $(0,\ldots,0,1)$.
Simultaneously scaling $w$ and the roots shows that the determinant
of the remaining $(n-1)$-square block has degree $n(n-1)$.
It must be a constant multiple of the discriminant.

Evaluate that constant at $P=w^n-a$, with $a$ indeterminate.
The block is diagonal, with
entries
\[
 -\frac{an(m+1)}{n(m+1)-j},\qquad 1\leq j\leq n-1.
\]
They follow by eliminating $w^{n+j-1}$ from
$w^{n+j-1}-aw^{j-1}$. The discriminant in the convention
of \eqref{cbc:hc:transition-determinant} is
\[
 \prod_{i<j}(z_i-z_j)^2
 =(-1)^{n(n-1)/2+n-1}n^na^{n-1}.
\]
Taking the product of the diagonal entries gives the formula.
For $n=1$ the transition is the identity.
\end{proof}

\begin{corollary}[Flatness and collision fibers]\label{cbc:hc:flat}
The $R_n$-module $H_n$ is flat. For every commutative
$R_n$-algebra $C$, its base change is the derivative quotient
formed from $C[w,\overline P^{-1}]/C[w]$.
Its geometric fiber is the residue space of the distinct points
remaining after collision.
\end{corollary}

\begin{proof}
The stage description expresses $H_n$ as a filtered colimit of
finite free modules. Each stage basis and its elimination remain
valid over $C$, because their pivots are nonzero integers in $k$.
Taking the colimit proves the base-change assertion. Over a
geometric point, powers of the specialized $P$ and of its reduced
polynomial give cofinal denominator systems. A zero-residue stage
class dies at the next transition. Every simple residue survives.
\end{proof}

\begin{example}[Two colliding points]\label{cbc:hc:binary}
Suppress the common translation coordinate and put
$R=k[u]$, $P=w(w-u)$. In the ordered basis
$a_m=P^{-m-1}$, $b_m=w^{2m+1}P^{-m-1}$, the residues at
$w=u$ are
\[
 \operatorname{Res}_{w=u}a_m\,dw
   =(-1)^m\binom{2m}{m}u^{-2m-1},\qquad
 \operatorname{Res}_{w=u}b_m\,dw=1.
\]
The total residues are respectively zero and one. Ratios of
successive first residues determine the transition on the
distinct locus. Since its entries are polynomials, this gives
the transition over $R$:
\begin{equation}\label{cbc:hc:binary-transition}
 T_m=\begin{pmatrix}
 -\dfrac{m+1}{2(2m+1)}u^2&0\\[3pt]0&1
 \end{pmatrix},\qquad
 H_2\simeq R[u^{-1}]\oplus R.
\end{equation}

The coordinates on the right are the residue at $w=u$ and
the total finite residue. For $m\geq1$, the differential
$dw/[w^m(w-u)]$ has coordinates $(u^{-m},0)$, while
$dw/(w-u)$ has coordinates $(1,1)$. Together with their
$R$-multiples these give the displayed module. The transition
formula proves injectivity of this residue description.
In residue coordinates the position connection is ordinary
differentiation in $u$: differentiating a moving simple pole
adds a double pole, whose residue is zero.

At $u=0$ the localized summand has zero fiber and the regular
summand remains. Thus flatness here does not imply finite
projectivity or constancy of fiber dimension.
\end{example}

\begin{proposition}\label{cbc:hc:no-section}
In Example~\ref{cbc:hc:binary}, the quotient $M_2\to H_2$
has no $R$-linear section. Hence the complex
$[M_2\xrightarrow{\partial_w}M_2]$ has no $R$-linear strong
deformation retract onto its cohomology on these modules.
\end{proposition}

\begin{proof}
The free module \eqref{cbc:hc:free-basis} satisfies
$\bigcap_{N\geq0}u^NM_2=0$. Every element in the image of an
$R$-linear map $R[u^{-1}]\to M_2$ lies in that intersection,
so the map is zero. A section would be injective on the
localized summand of \eqref{cbc:hc:binary-transition}.
A strong deformation retract would in particular supply
such a section.
\end{proof}

\section{The differential module on the collision diagonals}
\label{cbc:hc:diagonals}

Position derivatives retain information beyond a fiber's residue
space. Let
\[
 D_n=R_n\langle\partial_1,\ldots,\partial_n\rangle,
 \qquad [\partial_i,z_j]=\delta_{ij},\quad
 [\partial_i,\partial_j]=0,
\]
where $\partial_i$ differentiates a rational representative at
fixed $w$. For every nonempty $S\subseteq[n]$, define
\begin{equation}\label{cbc:hc:subset-generator}
 g_S=\left[\prod_{i\in S}(w-z_i)^{-1}\right]\in H_n.
\end{equation}
For a singleton write $g_i=g_{\{i\}}$.
Divided position derivatives recover arbitrary pole orders:
\begin{equation}\label{cbc:hc:divided-derivatives}
 \left[\prod_{i\in S}(w-z_i)^{-m_i}\right]
 =\prod_{i\in S}\frac{\partial_i^{m_i-1}}{(m_i-1)!}\,g_S,
 \qquad m_i\geq1.
\end{equation}
The finite list of $g_S$ can therefore describe an unbounded
range of pole orders.

Let $L_sM_n$ be the sum of the principal-part submodules whose
poles use at most $s$ labels, and let $F_sH_n$ be its image in
$H_n$. Set $L_0M_n=F_0H_n=0$. These filtrations are preserved
by position derivatives, and $\partial_w$ preserves $L_sM_n$.

For nonempty $S$, the diagonal $\Delta_S$ is given by
$z_i=z_j$ for $i,j\in S$. Denote by $\delta_S$ the cyclic
$D_n$-module generated by $\eta_S$, subject to
\begin{equation}\label{cbc:hc:delta-relations}
 (z_i-z_j)\eta_S=0\ (i,j\in S),\qquad
 \partial_k\eta_S=0\ (k\notin S),\qquad
 \sum_{i\in S}\partial_i\eta_S=0.
\end{equation}

It is the direct-image differential module of the trivial
connection on $\Delta_S$.

To see its normal directions explicitly,
choose $r=\max S$ and put $u_i=z_i-z_r$ for $i\in S\setminus\{r\}$.
Over the tangential polynomial ring it has the principal-part
basis $\prod_{i\ne r}u_i^{-b_i}$, $b_i\geq1$.
Multiplication by $u_i$ kills a term when $b_i=1$, and otherwise
lowers $b_i$. Differentiation sends $u_i^{-b_i}$ to
$-b_i u_i^{-b_i-1}$. For $|S|=1$, $\delta_S=R_n$.

Filter $D_n$ by differential order and write
$\xi_i=\operatorname{gr}(\partial_i)$.
A compatible exhaustive filtration on a module is good if
it is bounded below and its associated graded is finitely
generated over $R_n[\xi_1,\ldots,\xi_n]$.
The characteristic variety is the support of that associated
graded module in the cotangent bundle $T^*\mathbb A^n$.
A nonzero module is holonomic when that variety has dimension
$n$. Its holonomic length is its composition length.

\begin{theorem}[Diagonal constituents]\label{cbc:hc:filtration}
The pole-label filtration has differential-module quotients
\begin{equation}\label{cbc:hc:graded}
 F_sH_n/F_{s-1}H_n\simeq\bigoplus_{|S|=s}\delta_S,
 \qquad 1\leq s\leq n,
\end{equation}
with the image of $g_S$ as cyclic generator of its summand.
For $2\leq s\leq n$, the resulting extension of
$\bigoplus_{|S|=s}\delta_S$ by $F_{s-1}H_n$ does not split
even as a sequence of $R_n$-modules. The holonomic length of
$H_n$ is $2^n-1$. Over $k=\mathbb C$ it is regular holonomic.
\end{theorem}

\begin{proof}
The variables $w$ and $t_i=w-z_i$, $1\leq i\leq n$, are
independent coordinates on the auxiliary affine space.
A Laurent monomial determines its set of negative $t_i$
exponents. Modulo monomials with fewer negative labels, this
gives
\[
 \operatorname{gr}_s^L M_n
 =\bigoplus_{|S|=s}H^s_{(t_i:i\in S)}(R_n[w]).
\]
Here the local cohomology summand consists of negative powers
in all its indicated variables, with polynomial coefficients
in the remaining variables. This description is the top term
of the coordinate \v Cech complex and fixes its differential
module structure.

On the summand for $S$, use $t=w-z_r$, the normal coordinates
$u_i$, and the tangential coordinates $z_r,z_k$ for $k\notin S$.
The support ideal becomes $(t,u_i:i\ne r)$.
Differentiation in $w$ at fixed positions becomes $\partial_t$.
Thus the local cohomology module has basis
\[
 t^{-a}\prod_{i\in S\setminus\{r\}}u_i^{-b_i},
 \qquad a,b_i\geq1,
\]
over the tangential polynomial ring. The derivative in $t$ is
injective and its cokernel retains precisely $a=1$.

The induced position actions are \eqref{cbc:hc:delta-relations}.
The linear change of normal coordinates has determinant a
unit. Its fundamental fraction sends $g_S$ to a nonzero
multiple of the fundamental normal class. Normalize $\eta_S$
by this image.

Apply the kernel--cokernel sequence for $\partial_w$ to
$0\to L_{s-1}M_n\to L_sM_n\to\operatorname{gr}_s^LM_n\to0$.
The injectivity on every graded summand proves, by induction,
that the cokernels embed in $H_n$ and have the quotients
\eqref{cbc:hc:graded}.

By Corollary~\ref{cbc:hc:flat}, $H_n$ is torsion-free over the
domain $R_n$. For $s\geq2$, each element in a summand of
\eqref{cbc:hc:graded} is killed by a power of a normal equation.
A product of these equations kills an element of the finite
direct sum. Every $R_n$-linear map from that
quotient into $F_sH_n$ is therefore zero. It cannot be a section.

Each $\delta_S$ is simple. In the displayed basis, multiplication
by normal coordinates lowers the normal derivative degree,
and tangential differentiation lowers polynomial degree.
Choose a term of maximal total normal derivative degree and
multiply by its normal coordinate monomial. This kills all
other normal monomials of that degree and every lower one,
leaving a nonzero tangential polynomial times the fundamental
class. A derivative of maximal total tangential degree
extracts a nonzero scalar. Normal derivatives and tangential
multiplication then generate the whole module.

Its characteristic variety is the conormal
bundle of $\Delta_S$, of dimension $n$, so it is holonomic.
The finite filtration has one such constituent for each
nonempty subset, proving its length. Over $\mathbb C$ the
trivial diagonal connections are regular, and regular
holonomic modules are closed under extensions
\cite{HottaTakeuchiTanisaki2008}*{Definition~6.1.1}.
\end{proof}

The diagonal quotients are not flat over $R_n$. Consequently
their ordinary exterior powers cannot compute the associated
graded of exterior powers of $H_n$ without resolving the
intersections of their supports.

\begin{theorem}[Finite differential-module presentation]
\label{cbc:hc:presentation}
The generators \eqref{cbc:hc:subset-generator} present $H_n$
over $D_n$ with the relations
\begin{align}
 \partial_k g_S&=0 &&(k\notin S),\notag\\
 \sum_{i\in S}\partial_i g_S&=0,\notag\\
 (z_i-z_j)g_S&=g_{S\setminus\{j\}}-g_{S\setminus\{i\}}
       &&(i,j\in S,\ i\ne j).
 \label{cbc:hc:presentation-relations}
\end{align}
No empty-subset generator is required.
\end{theorem}

\begin{proof}
The first and third equations hold on rational representatives.
The second represents the $w$-derivative of the negative of
their product. Formula~\eqref{cbc:hc:divided-derivatives}
generates every pole multiplicity. Repeatedly writing
$w=(w-z_i)+z_i$ reduces any numerator to those fractions
and to polynomial terms, which vanish in $M_n$.
The presented module $Q$ therefore surjects onto $H_n$.

Filter $Q$ by the cardinalities of its generators. On its
level-$s$ quotient the right side of the last relation
vanishes. The remaining relations give surjections
\[
 \bigoplus_{|S|=s}\delta_S\longrightarrow
 \operatorname{gr}_sQ\longrightarrow\operatorname{gr}_s^FH_n.
\]
Their composite is the isomorphism of
Theorem~\ref{cbc:hc:filtration} with its chosen generators.
Both arrows are thus isomorphisms. Induction on the finite
filtration proves $Q=H_n$.
\end{proof}

\begin{example}[The triple attachment]\label{cbc:hc:triple}
For $a=z_1-z_3$ and $b=z_2-z_3$, the binary and ternary
generators satisfy
\begin{align*}
 (a-b)g_{12}&=g_1-g_2,& ag_{13}&=g_1-g_3,&
 bg_{23}&=g_2-g_3,\\
 ag_{123}&=g_{12}-g_{23},&
 bg_{123}&=g_{12}-g_{13}.&&
\end{align*}
On the distinct locus the residue vector of $g_{123}$ is
\[
 \left(\frac1{a(a-b)},\frac1{b(b-a)},\frac1{ab}\right).
\]
Its total residue is zero, but its image in
$H_3/F_2H_3=\delta_{123}$ is nonzero. The two attaching
equations specify how that constituent meets the three
binary diagonals.
\end{example}

\section{Collision and cut maps}
\label{cbc:hc:collision-cuts}

Let $\alpha:[n]\twoheadrightarrow[m]$ identify the positions
by $z_i=y_{\alpha(i)}$. The specialized coefficient has its
induced tangential connection: differentiation in $y_t$
acts by the sum of position derivatives in the fiber
$\alpha^{-1}(t)$. For nonempty $S$, set
$T=\alpha(S)$ and $\nu_t=|S\cap\alpha^{-1}(t)|$.

\begin{proposition}\label{cbc:hc:collisions}
The coefficient base-change isomorphism
$H_n\otimes_{R_n}k[y_1,\ldots,y_m]\simeq H_m$ sends
\begin{equation}\label{cbc:hc:collision-formula}
 g_S\longmapsto
 \left(\prod_{t\in T}
       \frac{\partial_{y_t}^{\nu_t-1}}{(\nu_t-1)!}\right)g_T.
\end{equation}
These maps are horizontal and compose under successive
surjections.
\end{proposition}

\begin{proof}
Both sides of \eqref{cbc:hc:collision-formula} are represented
by $\prod_{t\in T}(w-y_t)^{-\nu_t}$. Powers of the repeated
and reduced denominator products are cofinal, so substitution
identifies their localizations and principal parts.
Corollary~\ref{cbc:hc:flat} identifies their derivative
cokernels. Successive substitutions agree as rational
functions, and differentiation obeys the chain rule.
\end{proof}

If all labels of $S$ collapse to one point and $|S|>1$, its
image is zero, since
\[
 (w-y)^{-|S|}
 =\partial_w\left(-\frac{(w-y)^{1-|S|}}{|S|-1}\right).
\]
The proposition concerns coefficient base change with the
specified tangential action. An extraordinary inverse image
of differential modules also involves its density and shift
conventions.

A consecutive cut separates groups of labels while allowing
collisions inside each group. More generally, let
$\mathcal B=(B_1|\cdots|B_r)$ be an ordered partition of
$[n]$ into nonempty blocks. Let $U_{\mathcal B}$ be the open
where $z_i-z_j$ is invertible whenever $i$ and $j$ lie in
different blocks. Write $R_{\mathcal B}$ for its coordinate
ring and $M_{B_a},H_{B_a}$ for the coefficients using just
the labels in $B_a$, extended to $R_{\mathcal B}$.

\begin{proposition}[Block decomposition]\label{cbc:hc:cuts}
There are canonical horizontal isomorphisms
\begin{equation}\label{cbc:hc:cut-decomposition}
 M_n|_{U_{\mathcal B}}\simeq\bigoplus_aM_{B_a},\qquad
 H_n|_{U_{\mathcal B}}\simeq\bigoplus_aH_{B_a}.
\end{equation}
They commute with ordered refinements and with collisions
inside the blocks.
\end{proposition}

\begin{proof}
Put $P_a=\prod_{i\in B_a}(w-z_i)$. The resultants of
different $P_a$ are units on $U_{\mathcal B}$, so their
powers are pairwise coprime. The polynomial Chinese
remainder theorem gives a unique proper partial-fraction
decomposition by these denominators. This proves the first
isomorphism.

Differentiation in $w$ preserves every summand,
giving the second. Position differentiation also preserves
each group's allowed poles, so the maps are horizontal.
Uniqueness makes successive refinements agree. After an
internal collision the resultants between groups remain
units. The same decomposition commutes with substitution.
\end{proof}

\section{Free comparisons and completed mapping complexes}
\label{cbc:hc:free-comparison}

The absence of a section in Proposition~\ref{cbc:hc:no-section}
does not prevent comparison by free resolutions. The finite
stages supply both the resolution and its chain map.
Let $E=\bigoplus_{m\geq0}F_m$, and let $\sigma_T:E\to E$
send the $m$-th summand to the next by $T_m$.

\begin{proposition}\label{cbc:hc:telescope}
The sequence
\begin{equation}\label{cbc:hc:telescope-sequence}
 0\longrightarrow E\xrightarrow{1-\sigma_T}E
 \longrightarrow H_n\longrightarrow0
\end{equation}
is a free $R_n$-resolution. For an $R_n$-module $N$,
it computes derived Hom as
\begin{equation}\label{cbc:hc:mapping-complex}
 \begin{aligned}
 R\operatorname{Hom}_{R_n}(H_n,N)
 &=\left[\prod_m\operatorname{Hom}_{R_n}(F_m,N)
 \xrightarrow{\delta_T}
 \prod_m\operatorname{Hom}_{R_n}(F_m,N)\right],
 \\
 (\delta_Tf)_m&=f_m-f_{m+1}T_m.
 \end{aligned}
\end{equation}
in degrees zero and one.
\end{proposition}

\begin{proof}
The first nonzero stage of a finite-support vector survives
under $1-\sigma_T$, proving injectivity. Its cokernel is
the defining colimit. Both copies of $E$ are free.
Applying Hom to these direct sums gives the products in
\eqref{cbc:hc:mapping-complex} and the stated difference map.
\end{proof}

The elimination in Theorem~\ref{cbc:hc:stages} records, along
with each matrix column, a primitive $p_{m,j}\in M^{(m+1)}$:
\begin{equation}\label{cbc:hc:transition-primitive}
 c_{m,j}-\sum_i(T_m)_{ij}c_{m+1,i}=\partial_wp_{m,j}.
\end{equation}
Each eliminated numerator contributes its coefficient times
the rational function whose derivative was eliminated.

Place the two copies of $E$ in degrees zero and one,
with corresponding generators $u_{m,j}$ and $v_{m,j}$.
The differential is
$du_{m,j}=v_{m,j}-\sum_i(T_m)_{ij}v_{m+1,i}$.

\begin{proposition}\label{cbc:hc:chain-comparison}
The maps
\begin{equation}\label{cbc:hc:comparison-formula}
 u_{m,j}\longmapsto-\kappa^{-1}x(p_{m,j}),\qquad
 v_{m,j}\longmapsto\gamma(c_{m,j})
\end{equation}
give an $R_n$-linear homotopy equivalence from this
telescope to $K_{\kappa,n}$.
\end{proposition}

\begin{proof}
Equation~\eqref{cbc:hc:transition-primitive} is exactly the
chain-map identity. Both degree-zero cohomologies vanish,
and the map is the identity on degree-one cohomology $H_n$.
Both complexes are bounded and free over $R_n$.

An inverse can be constructed on their free bases. Reduce
each degree-one basis vector of $M_n$ at a finite stage
and lift its class to the corresponding $F_m\subset E$.

The image of a boundary under this lift lies in
$\ker(E\to H_n)=\operatorname{im}(1-\sigma_T)$. Its unique
preimage defines the degree-zero map. To find the preimage,
advance its finitely many components to a common stage and
telescope the intervening transitions.
Differences between
either composite and the identity induce zero on $H_n$.
Lift those differences through the injective boundary maps
in the same way. These lifts are the chain homotopies.
\end{proof}

For $P=w(w-u)$ the first two primitives are
\begin{equation}\label{cbc:hc:binary-primitives}
 p_{0,\mathrm{low}}=-\frac{w-u/2}{P},\qquad
 p_{0,\mathrm{top}}=\frac{u}{w-u}.
\end{equation}
Their derivatives are $P^{-1}+u^2/(2P^2)$ and
$(w-u)^{-1}-w/(w-u)^2$. Even the identity transition on
the top cohomology coordinate therefore has a nonzero
chain correction. The displayed homotopy equivalence is
$R_n$-linear. A connection on a chosen resolution requires
its own lifts of position differentiation.

\begin{proposition}[The effect of a connection]\label{cbc:hc:reverse-hom}
For $R=k[u]$ and $L=R[u^{-1}]$,
\[
 \operatorname{Hom}_R(L,R)=0,\qquad
 \operatorname{Ext}^1_R(L,R)=k[[u]]/R,
 \qquad \operatorname{Ext}^{q}_R(L,R)=0\quad(q>1).
\]
Give $L$ and $R$ their ordinary derivative connections, and
put $D=k\langle u,\partial\rangle/([\partial,u]-1)$.
Then $R\operatorname{Hom}_D(L,R)=0$.
With target $L$, the groups $\operatorname{Hom}_D(L,L)$
and $\operatorname{Ext}^1_D(L,L)$ are each $k$, and the
higher groups vanish.
\end{proposition}

\begin{proof}
Present $L$ by the telescope with generators $u^{-2m}$
and transition $u^2$. Its Hom differential is
$(f_m)\mapsto(f_m-u^2f_{m+1})$ on $\prod_{m\geq0}R$.
The following series converges in the $u$-adic topology.
The map
\[
 (a_m)\longmapsto\sum_{m\geq0}u^{2m}a_m\pmod R
\]
takes boundaries to polynomials. If the sum is the
polynomial $p$, then
\[
 f_m=u^{-2m}\left(p-\sum_{j<m}u^{2j}a_j\right)\in R
\]
satisfies $a_m=f_m-u^2f_{m+1}$.

Grouping the even and odd
digits of any power series proves surjectivity.
Infinite $u$-divisibility makes the degree-zero kernel
zero, and the resolution has length one.

For the differential-module calculation, the assignment
$e\mapsto u^{-1}$ gives
\begin{equation}\label{cbc:hc:localization-presentation}
 L\simeq D/D(u\partial+1).
\end{equation}
Indeed, commuting $u$ past powers of $\partial$ reduces the
quotient to the spanning set $u^je$ for $j\geq0$ and
$\partial^je$ for $j\geq1$. Their images are respectively
$u^{j-1}$ and $(-1)^jj!u^{-j-1}$, distinct Laurent
monomials, so they are independent.

The free resolution
given by right multiplication by $u\partial+1$ is
injective, by the order filtration on the Weyl algebra.
Applying Hom gives $u\partial+1$ on the target. On $R$
its eigenvalue at $u^j$ is $j+1$, always invertible.
On $L$ its kernel and cokernel are both the line represented
by $u^{-1}$.
\end{proof}

For example, the constant sequence $a_m=1$ represents the
nonzero ordinary extension class of $(1-u^2)^{-1}$ modulo
$R$. That extension has generators $e,f_0,f_1,\ldots$ and
relations $f_m-u^2f_{m+1}=e$. It cannot carry the prescribed
connections as an extension of $L$ by $R$, since the
corresponding differential-module extension group vanishes.

\section{A telescope with all spatial derivatives}
\label{cbc:hc:spatial}

Position differentiation can raise the denominator stage. To keep
it on a free resolution, the relation generators must move
as well. This can be done before taking the auxiliary
derivative quotient.

Let $d\geq0$, let $B=k[s_1,\ldots,s_d]$, and let $f\in B[w]$ be monic
of positive $w$-degree. Define the relative Weyl algebra
\[
 D_B=B\langle w,\partial_w\rangle/([\partial_w,w]-1),
\]
in which the parameters are central. Write $f_a=\partial_{s_a}f$,
$f_w=\partial_wf$, and $f_{wa}=\partial_w\partial_{s_a}f$.
Put
\begin{equation}\label{cbc:hc:weyl-operators}
 A_m=f\partial_w+mf_w,\qquad
 A_{m,a}=f_a\partial_w+mf_{wa},\qquad m\geq1.
\end{equation}
These are operators: $\partial_wg=g\partial_w+g_w$.
The identity $A_mf=fA_{m+1}$ gives transition maps, by
right multiplication by $f$, between the free complexes
$[D_B\xrightarrow{r_{A_m}}D_B]$ in degrees $-1,0$.

Their direct-sum telescope $\mathcal T_f$ has generators
$e_m$ in degree zero, $r_m,q_m$ in degree $-1$, and $h_m$
in degree $-2$. Define its $D_B$-linear differential by
\begin{equation}\label{cbc:hc:spatial-differential}
 \begin{aligned}
 de_m&=0,& dr_m&=A_me_m,\\
 dq_m&=e_m-fe_{m+1},&
 dh_m&=r_m-fr_{m+1}-A_mq_m.
 \end{aligned}
\end{equation}
Every vector has finite stage support. The only potentially
nonzero square is
$d^2h_m=(A_mf-fA_{m+1})e_{m+1}=0$.

\begin{lemma}\label{cbc:hc:localization-telescope}
The augmentation $e_m\mapsto f^{-m}$, with negative-degree
generators sent to zero, makes $\mathcal T_f$ a free
$D_B$-resolution of $B[w,f^{-1}]$. After every scalar
extension $B\to C$ to a commutative $k$-algebra it remains
a resolution of $C[w,\overline f^{-1}]$ and computes
derived scalar extension.
\end{lemma}

\begin{proof}
Right multiplication by $A_m$ is injective: its leading
coefficient in differential order is multiplication by
the monic polynomial $f$. This remains a non-zero-divisor
over $C$. Thus each two-term complex resolves
$P_m=D_B/D_BA_m$.

In their colimit the transition and the next relation give
\[
 e_m=fe_{m+1},\qquad
 \partial_we_m=-mf_we_{m+1}.
\]
Repeatedly applying the second equation reduces every
operator applied to $e_m$ to a polynomial times a later
generator. The fraction map has inverse $p/f^m\mapsto pe_m$.
The first equation makes it well-defined and the second
makes it horizontal. Hence $\operatorname*{colim}P_m=B[w,f^{-1}]$.

The map $1-\mathrm{shift}$ on the direct sum of the stage
complexes is injective at the first nonzero stage. Its cone
therefore resolves the filtered colimit, and its differential
is \eqref{cbc:hc:spatial-differential} with the displayed
choice of sign on $h_m$.

Its terms are free over $D_B$
and over $B$. The same arguments after scalar extension
prove the last assertion.
\end{proof}

\begin{example}[A finite stage at collision]\label{cbc:hc:spurious-delta}
Take $f=w(w-s)$ over $k[s]$ and specialize to $s=0$.
Write $D=k\langle w,\partial_w\rangle/([\partial_w,w]-1)$.
Then $A_m=wB_m$, where $B_m=w\partial_w+2m$.
In the specialized $P_m$, put $g_m=B_me_m$.
The identities
\[
 wg_m=0,\qquad
 (B_m-(2m-1))B_m=\partial_w A_m
\]
give $B_mg_m=(2m-1)g_m$.

Consequently $\ell_m=e_m-g_m/(2m-1)$ is killed by $B_m$,
and
\begin{equation}\label{cbc:hc:finite-collision}
 P_m|_{s=0}\simeq k[w,w^{-1}]\oplus D/Dw.
\end{equation}

For an explicit inverse, send $e_m$ to
$(w^{-2m},\eta/(2m-1))$, where $w\eta=0$.
The operator $A_m$ kills this pair. The element $g_m$ maps to $(0,\eta)$
and $\ell_m$ to $(w^{-2m},0)$. These maps verify both
summands of \eqref{cbc:hc:finite-collision}.

The extra delta is removed by the next denominator transition:
\[
 g_m\longmapsto B_mw^2e_{m+1}
 =w^2B_{m+1}e_{m+1}=w^2g_{m+1}=0.
\]
The whole telescope therefore has the correct collision
fiber, although a single cyclic stage does not.
\end{example}

Extend $\partial_{s_a}$ to $D_B$ by fixing $w,\partial_w$.
Define $\nabla_a(Pv)=(\partial_{s_a}P)v+P\nabla_av$ and
\begin{equation}\label{cbc:hc:spatial-connection}
 \begin{aligned}
 \nabla_ae_m&=-mf_ae_{m+1},\\
 \nabla_ar_m&=-(m-1)f_ar_{m+1}+A_{m,a}q_m,\\
 \nabla_aq_m&=-(m+1)f_aq_{m+1},\\
 \nabla_ah_m&=-mf_ah_{m+1}.
 \end{aligned}
\end{equation}
The term $A_{m,a}q_m$ compares the derivative of a relation
with its image at the next stage.

\begin{theorem}[Flat spatial connection]\label{cbc:hc:spatial-flatness}
The operators \eqref{cbc:hc:spatial-connection} satisfy
\[
 [d,\nabla_a]=0,\qquad [\nabla_a,\nabla_b]=0,\qquad
 [\nabla_a,s_b]=\delta_{ab},\qquad
 [\nabla_a,w]=[\nabla_a,\partial_w]=0.
\]
They make $\mathcal T_f$ a complex over the full Weyl algebra
$\mathscr D=D_B\langle\partial_{s_1},\ldots,\partial_{s_d}\rangle$,
with horizontal augmentation to $B[w,f^{-1}]$ with its
ordinary derivatives. No discriminant is inverted.
\end{theorem}

\begin{proof}
The scalar commutators follow from the coefficient Leibniz
rule. For the chain equation on $r_m$, normal ordering gives
\begin{equation}\label{cbc:hc:mixed-identity}
 A_{m,a}f+(m-1)f_aA_{m+1}=mA_mf_a,
\end{equation}
both sides equal
$mff_a\partial_w+mff_{wa}+m^2f_wf_a$.
On $q_m$, both orders of differentiation give
$-(m+1)f_ae_{m+1}+(m+1)ff_ae_{m+2}$.
On $h_m$, both give
\[
 -mf_ar_{m+1}+mff_ar_{m+2}+mf_aA_{m+1}q_{m+1}.
\]
The last equality follows by differentiating
$A_mf=fA_{m+1}$ and using \eqref{cbc:hc:mixed-identity}.
The chain equation on $e_m$ is immediate.

For curvature, mixed derivatives and scalar products cancel
on $e_m,q_m,h_m$. On $r_m$, the $q_m$ coefficient cancels
by equality of mixed derivatives of $A_m$. The remaining
$q_{m+1}$ coefficient is
\begin{align*}
 &(m-1)(f_aA_{m+1,b}-f_bA_{m+1,a})\\
 &\hspace{8mm}+(m+1)(A_{m,a}f_b-A_{m,b}f_a).
\end{align*}
Writing $W=f_af_{wb}-f_bf_{wa}$, the first line is
$(m^2-1)W$ and the second is $(m+1)(1-m)W$.
Their sum is zero. Finally,
$\partial_{s_a}f^{-m}=-mf_af^{-m-1}$ proves horizontality
of the augmentation.
\end{proof}

For $f=w^3+aw+b$, the correction operators are
$A_{m,a}=w\partial_w+m$ and $A_{m,b}=\partial_w$.
These formulas remain defined on $4a^3+27b^2=0$.
Deleting the mixed term from $\nabla_br_1$ would give
the nonzero chain defect
$-\partial_we_1+\partial_wf\,e_2$, where $\partial_wf$
denotes operator composition. The correction
$d(\partial_wq_1)$ cancels it.

To obtain free terms over the full Weyl algebra, let
$v_1,\ldots,v_d$ be the standard basis of $k^d$, and form
\[
 \mathcal P_f^j
 =\bigoplus_{q-p=j}
   \mathscr D\otimes_{D_B}\mathcal T_f^q\otimes_k\Lambda^pk^d,
 \qquad 0\leq p\leq d.
\]
On the induced module set
$L_a(P\otimes u)=P\partial_{s_a}\otimes u-P\otimes\nabla_au$.
The connection's Leibniz rule makes $L_a$ balanced over
$D_B$. Let $\delta$ be their alternating sum deleting the
corresponding wedge factor, and define
\begin{equation}\label{cbc:hc:spencer-differential}
 d_{\mathcal P}(P\otimes u\otimes\omega)
 =P\otimes du\otimes\omega
   +(-1)^q\delta(P\otimes u\otimes\omega),
 \qquad u\in\mathcal T_f^q.
\end{equation}
Flatness makes the $L_a$ commute, and the chain equation
makes them commute with the internal differential.
The total sign in \eqref{cbc:hc:spencer-differential}
therefore gives $d_{\mathcal P}^2=0$.

\begin{proposition}\label{cbc:hc:spencer-resolution}
The action augmentation
$\mathcal P_f\to\mathcal T_f\to B[w,f^{-1}]$ is a free
$\mathscr D$-resolution, bounded in degrees $-d-2,\ldots,0$.
\end{proposition}

\begin{proof}
Fix $q$ and normally order the parameter derivatives.
The induced module becomes their polynomial module with
coefficients in $\mathcal T_f^q$. Filter the exterior-degree
$p$ term by parameter-derivative order plus $p$.
The leading part of $L_a$ is multiplication by its independent
polynomial variable. The associated augmented row is thus
the Koszul resolution of $\mathcal T_f^q$ for those variables.
It is exact.

Lifting a leading boundary lowers the filtration.
Every vector has finite order, so this procedure terminates.
The augmented row resolves $\mathcal T_f^q$.

The internal degrees are bounded, and totalization gives
a quasi-isomorphism to $\mathcal T_f$.
Lemma~\ref{cbc:hc:localization-telescope} identifies its
cohomology. Each induced term is free over $\mathscr D$
because $\mathcal T_f^q$ is free over $D_B$.
\end{proof}

For $B=R_n$ and $f=P$, this resolves the auxiliary localization
in \eqref{cbc:hc:poles}. The primitive coefficient $H_n$ is
obtained by first quotienting by $R_n[w]$ and then taking
the $\partial_w$-cokernel. It is a different module.
Applying Hom to the free spatial resolution takes products
over all its stage generators. Those products are part of
the mapping complex and cannot be replaced by a finite
denominator cutoff.

\section{The polynomial Heisenberg complex}
\label{cbc:hc:polynomial}

The primitive differential \eqref{cbc:hc:linear-complex}
extends as a derivation to
\begin{equation}\label{cbc:hc:polynomial-complex}
 Z_{\kappa,n}=\operatorname{Sym}_{R_n}(xM_n)
      \otimes_{R_n}\Lambda_{R_n}(\gamma M_n),\qquad
 dx(f)=-\kappa\gamma(\partial_wf),\quad d\gamma(f)=0.
\end{equation}
Here $x(f)$ has degree zero and $\gamma(f)$ degree one.
The product is the auxiliary graded-commutative polynomial
product, and each element has finite polynomial degree.

\begin{theorem}\label{cbc:hc:polynomial-cohomology}
The primitive projection extends to a quasi-isomorphism
of differential graded $R_n$-algebras
\begin{equation}\label{cbc:hc:cohomology-map}
 \rho_n:Z_{\kappa,n}\longrightarrow\Lambda_{R_n}(H_n),
 \qquad x(f)\longmapsto0,\quad\gamma(f)\longmapsto[f],
\end{equation}
where $H_n$ has degree one and zero differential.
It preserves the unit, the position connection, coefficient
collision maps, and the block decompositions of
Proposition~\ref{cbc:hc:cuts}.
\end{theorem}

\begin{proof}
The primitive map is a quasi-isomorphism between bounded
flat complexes by Theorem~\ref{cbc:hc:stages} and
Corollary~\ref{cbc:hc:flat}. Tensor powers preserve this
quasi-isomorphism. In characteristic zero, each graded
symmetric power is the image of its averaging idempotent
on a tensor power. It therefore preserves the map as well.
Their direct sum is \eqref{cbc:hc:cohomology-map}.

The generator formulas commute with differentiation,
substitution, and the direct-sum decomposition of principal
parts. They also preserve the multiplicative unit.
\end{proof}

Let $V$ be finite-dimensional with a nondegenerate level
form, and identify $V$ with its dual by that form.
Apply the same argument to the corresponding copies of
$V\otimes M_n$. The resulting cohomology is
$\Lambda_{R_n}(V\otimes H_n)$.

Over the normal-coordinate ring $R=k[u]$ of
Example~\ref{cbc:hc:binary}, the rank-one polynomial complex
has cohomology
\[
 H^0(Z_{\kappa,2})=R,\qquad
 H^1(Z_{\kappa,2})=R[u^{-1}]\oplus R,\qquad
 H^2(Z_{\kappa,2})=R[u^{-1}],
\]
with no higher cohomology. In particular
$\Lambda_R^2R[u^{-1}]=0$, since every pair of elements
lies in a common rank-one free submodule.
Under a cut, exterior algebra converts
\eqref{cbc:hc:cut-decomposition} into a tensor product.
Comparison with the transverse product requires its
contraction maps. The polynomial projection alone does
not specify those operations.

\section{The binary holonomic dual}
\label{cbc:hc:dual}

Duality must transpose the attachment, including its normal
derivative. In one variable let $D$ and $L$ be as in
Proposition~\ref{cbc:hc:reverse-hom}. Trivialize the canonical
line by $du$ and use the adjoint anti-involution
$u^t=u$, $\partial^t=-\partial$.
For a holonomic left module $N$, the convention is
\[
 \mathbb D N={}^{t}R\operatorname{Hom}_D(N,D)[1],
\]
where a right module becomes a left module by the adjoint.
This includes the one-dimensional duality shift
\cite{HottaTakeuchiTanisaki2008}*{Definition~2.6.1 and Corollary~2.6.8}.

\begin{proposition}\label{cbc:hc:binary-dual}
There is an isomorphism
\begin{equation}\label{cbc:hc:dual-module}
 \mathbb D L\simeq D/D(u\partial).
\end{equation}
If $e$ is its cyclic generator, then $p_1=\partial e\ne0$,
and there is a nonsplit exact sequence of differential modules
\begin{equation}\label{cbc:hc:dual-extension}
 0\longrightarrow D/Du\xrightarrow{\ 1\mapsto p_1\ }
 D/D(u\partial)\xrightarrow{\ e\mapsto1\ }k[u]
 \longrightarrow0.
\end{equation}
This sequence splits as a sequence of $k[u]$-modules.
\end{proposition}

\begin{proof}
Apply $\operatorname{Hom}_D(-,D)$ to the free resolution
of \eqref{cbc:hc:localization-presentation}. Its only
cohomology is in degree one. The adjoint of the relation is
\[
 (u\partial+1)^t=-\partial u+1=-u\partial.
\]
The duality shift puts this cokernel in degree zero,
giving \eqref{cbc:hc:dual-module}.

A normal spanning set for the quotient is $u^me$, $m\geq0$,
and $p_j=\partial^je$, $j\geq1$. Their actions are
\begin{align*}
 up_1&=0,& up_j&=-(j-1)p_{j-1}\quad(j>1),&
 \partial p_j&=p_{j+1},\\
 \partial e&=p_1,&
 \partial(u^me)&=m u^{m-1}e\quad(m\geq1).&&
\end{align*}
These rules satisfy $[\partial,u]=1$ on every indicated
vector and $u\partial e=0$. Conversely they define a
module on the free vector space with that basis, generated
by $e$. The normal spanning set is therefore independent.

The span of the $p_j$ is $D/Du$, and the quotient on the
$u^me$ is $k[u]$, proving exactness.

Every lift of $1\in k[u]$ is $e+\sum_{j=1}^Na_jp_j$.
Its derivative has nonzero $p_1$ coefficient, so no lift is
horizontal and no $D$-linear section exists. The assignment
$u^m\mapsto u^me$ is, however, a $k[u]$-linear section.
\end{proof}

The original localization sequence is
$0\to k[u]\to L\to L/k[u]\to0$.
The adjoints of their relations $\partial$ and $u$ show
that the regular and delta constituents are self-dual.
Sequence~\eqref{cbc:hc:dual-extension} reverses their order and
retains the nonzero derivative $\partial e=p_1$.
Tensoring with the tangential regular connection gives the
corresponding extension with the common translation coordinate.
By \eqref{cbc:hc:binary-transition}, the complete primitive
dual over the normal-coordinate ring is
\[
 \mathbb D H_2\simeq k[u]\oplus D/D(u\partial).
\]

This calculation fixes the extension in the primitive
holonomic category. A duality for the full field coefficient
must additionally specify its completion and its operations.
